\chapter{示性类与Chern--Weil理论}
\label{ch:18-characteristic}
\ChapterPlan{周第8章（本地页码待核对）；Meinrenken、Freed讲义及Milnor原论文补充}{正文证明初稿；8题详解与3幅原创图，尚未独立数学审阅}

第16章的曲率本来依赖联络，第17章却已看到某些曲率积分不能通过换联络消去。本章解释这件事的机制：先对曲率代入共轭不变多项式，再用Bianchi把它变成闭形式；改变联络所造成的差恰为一个全局外微分。所得不变量首先生活在实de Rham上同调中，不能仅靠曲率公式便获得整系数中的挠信息或完整的丛分类。

流形仍有限维、Hausdorff、第二可数、光滑、无边界；涉及边界时另行声明。主丛用右作用，曲率为$\Omega=\dd\omega+[\omega,\omega]/2$，向量丛中为$F=\dd A+A\wedge A$。复丛沿用第17章归一化$\det(I+iF/(2\pi))$，因此Hopf相伴重言线丛的一阶Chern周期为$-1$。本文所称Chern、Pontryagin、Euler\emph{类}若未加上标，均指本章构造的实de Rham类；与拓扑整类的比较单列外引。

\Needspace{16\baselineskip}
\section{不变多项式与Chern--Weil构造}
\subsection{极化、次数与全群不变性}
设$G$为李群，李代数为$\mathfrak g$。齐次$k$次多项式$p:\mathfrak g\to\mathbb K$与对称$k$线性型$P$对应，约定
\[
 p(X)=P(X,\ldots,X),\qquad
 P(X_1,\ldots,X_k)=\left.\frac1{k!}
 \frac{\partial^k}{\partial t_1\cdots\partial t_k}
 p\left(\sum_jt_jX_j\right)\right|_{t=0}.
\]
将齐次多项式按单项式展开可验证此极化公式，特别是$P(X,\ldots,X)=p(X)$；$k!$不能在代入曲率后再重复加入。$k=0$为常数。要求
\[
 p(\operatorname{Ad}_gX)=p(X)\quad(g\in G),
\]
等价于$P$对所有输入同时施加$\operatorname{Ad}_g$后不变。若$G$不连通，仅有李代数的无穷小不变性还不足以代替此条件；Euler的Pfaffian正好会区分$\mathrm{SO}(2r)$和$\mathrm O(2r)$。

取基$e_a$，对$\mathfrak g$值形式$\beta_j=\sum_a\beta_j^ae_a$定义
\begin{equation}\label{eq:polarized-forms}
 P(\beta_1,\ldots,\beta_k)=
 \sum_{a_1,\ldots,a_k}P(e_{a_1},\ldots,e_{a_k})
 \beta_1^{a_1}\wedge\cdots\wedge\beta_k^{a_k}.
\end{equation}
换基矩阵在系数及输入中互相抵消，所以定义不依赖基。交换两个次数$r,s$的输入，符号为$(-1)^{rs}$；曲率输入均为二次，彼此交换没有负号。多项式次数$k$由此变成形式次数$2k$。

\subsection{不变性怎样消去联络项}
在底空间的局部截面中记势为$A$，协变外微分$D_A\beta=\dd\beta+[A,\beta]$，括号楔积按第16章约定。

\Needspace{15\baselineskip}
\begin{proposition}\label{prop:invariant-derivative}
对齐次形式$\beta_j$，有
\begin{equation}\label{eq:invariant-derivative}
 \dd P(\beta_1,\ldots,\beta_k)
 =\sum_j(-1)^{\deg\beta_1+\cdots+\deg\beta_{j-1}}
 P(\beta_1,\ldots,D_A\beta_j,\ldots,\beta_k).
\end{equation}
\end{proposition}
\begin{proof}
从全群不变性对$g=\exp(t\xi)$求导，得到
\[
 \sum_jP(X_1,\ldots,[\xi,X_j],\ldots,X_k)=0.
\]
把$A=\sum_aA^ae_a$及各$\beta_j$按基展开。右端中含$\dd\beta_j$的项就是外微分乘法式。其余含$[A,\beta_j]$的项中，将一次形式$A^a$移到全部$\beta$前面时，恰再产生$(-1)^{\deg\beta_1+\cdots+\deg\beta_{j-1}}$，与原符号相消；剩下的系数按上一条无穷小恒等式求和为零。这说明消掉的是整组联络项，而不是假设$A$在某个开集为零。
\end{proof}

\Needspace{18\baselineskip}
\begin{theorem}[Chern--Weil闭形式]\label{thm:chern-weil-closed}
设$\mathcal P\to M$为主$G$丛，$\omega$为主联络。对$\operatorname{Ad}_G$不变的齐次$k$次$p$，总空间的$P(\Omega,\ldots,\Omega)$唯一写成$\pi^*\operatorname{cw}_p(\omega)$，其中$\operatorname{cw}_p(\omega)$是底空间的闭$2k$形式。该构造对主丛同构和拉回自然，并保持多项式的和与积。
\end{theorem}
\begin{proof}
第16章曲率水平且$\operatorname{Ad}$等变，故代入$P$后水平性仍在，不变性又使其对右群作用不变。沿局部截面拉回得到$P(F_i,\ldots,F_i)$；不同截面的曲率同时共轭，故这些式子相等并拼成底空间形式。$\pi$为满射下沉映射，其拉回对形式单射，保证唯一性。

局部用\eqref{eq:invariant-derivative}及Bianchi $D_AF_i=0$，得到
\[
 \dd P(F_i,\ldots,F_i)=kP(D_AF_i,F_i,\ldots,F_i)=0.
\]
丛映射拉回曲率，且拉回保持楔积，所以自然性成立。把两个多项式极化后的曲率代入展开，因所有曲率分量为偶次，结果等于普通多项式乘积的代入；故$\operatorname{cw}_{pq}=\operatorname{cw}_p\wedge\operatorname{cw}_q$，和的性质由线性性得到。
\end{proof}
第18.3节将不使用任何整类比较，直接证明$[\operatorname{cw}_p(\omega)]$与联络无关。届时得到次数加倍的代数同态
\[
 (S\mathfrak g^*)^G\longrightarrow H^{\mathrm{even}}_{\mathrm{dR}}(M;\mathbb K).
\]
本章不证明它对某个分类空间是同构，也不声称在一个固定$M$上能检出全部丛。

\begin{figure}[htbp]
\centering\DGDiagram{characteristic-descent}

\caption{曲率先经不变多项式变成水平不变形式，再下降到基底。闭性使用Bianchi；从闭形式到联络无关的类还需第18.3节的传递公式。}
\label{fig:characteristic-descent}
\end{figure}

\subsection{矩阵值形式的迹}
对$r$次和$s$次矩阵值形式$\alpha,\beta$，按矩阵指标求和并交换标量形式可得
\[
 \operatorname{tr}(\alpha\wedge\beta)=(-1)^{rs}\operatorname{tr}(\beta\wedge\alpha),
 \quad \dd\operatorname{tr}\alpha=\operatorname{tr}(D_A\alpha).
\]
第二式中的$D_A\alpha=\dd\alpha+A\wedge\alpha-(-1)^r\alpha\wedge A$，两联络项的迹由第一式抵消。因此$\dd\operatorname{tr}F^k=0$。这里使用\emph{带次数}的循环律，不是对所有形式都无符号循环；例如$\operatorname{tr}A^4=0$，因为把第一个一次$A$移过其余三个后得到负的自身。


\Needspace{16\baselineskip}
\section{Chern类、Pontryagin类与Euler类}
\subsection{Chern形式及其实性}
设$E$为秩$r$的光滑复向量丛，先选Hermite内积及保持它的联络。在酉标架中曲率矩阵满足$F^*=-F$。定义总Chern形式为
\begin{equation}\label{eq:total-chern}
 c(\nabla;t)=\det\left(I+t\frac{iF}{2\pi}\right)
 =\sum_{k=0}^r t^k c_k(\nabla),\qquad c_0=1.
\end{equation}
$t$在这里是记账用的形式参数，每个$c_k$为实$2k$形式。行列式在偶形式组成的交换代数中计算；若把一次联络矩阵$A$直接放入同一行列式，交换性依据就丢失了。

\Needspace{17\baselineskip}
\begin{proposition}\label{prop:chern-properties}
各$c_k(\nabla)$闭，换标架不变，且对拉回自然。其上同调类$c_k(E)$与联络及Hermite度量无关，并有
\begin{gather*}
 c(E\oplus F)=c(E)c(F),\qquad c_1(E)=c_1(\det E),\\
 c_k(E^*)=c_k(\bar E)=(-1)^kc_k(E).
\end{gather*}
前述等式是类的等式；选直和、对偶、行列式诱导联络时才直接成为相应形式的等式。
\end{proposition}
\begin{proof}
行列式的各系数为$\mathrm{GL}(r,\mathbb C)$共轭不变多项式，由定理\ref{thm:chern-weil-closed}给闭性、拼接和自然性。酉标架中$iF$满足共轭转置不变，行列式又不因转置改变，故所有系数为实形式。

联络无关性由定理\ref{thm:cw-transgression}给出，其证明只用不变性及Bianchi，不预用本命题。即使两端联络保持的是不同Hermite内积，也可在全体复线性联络的仿射空间中连接；中间形式可为复值，但两个实端点之差等于传递形式\emph{实部}的外微分，故给同一实类。

直和联络曲率为分块对角，行列式乘法即给Whitney公式。对偶联络曲率$-F^{\mathsf T}$把每个$k$次系数乘$(-1)^k$。共轭丛的曲率为$\bar F=-F^{\mathsf T}$（酉标架中），得到同一结论。行列式线丛的曲率为$\operatorname{tr}F$，其一阶形式与原丛的一阶形式相同。任意其他联络再由联络无关性比较。
\end{proof}
若$E$有无零光滑截面，归一化后由第16章分裂为$\mathbb C\oplus E'$，$E'$秩为$r-1$。Whitney公式及秩截断给$c_r(E)=0$；这仅是必要条件。平凡丛上仍可有非零的\emph{恰当}Chern形式，例如第17章平凡线丛取$H=e^{|z|^2}$，其曲率为$-\dd z\wedge\dd\bar z$，对应形式非零而类为零。

\subsection{Chern特征与低阶展开}
记$B=iF/(2\pi)$，定义
\begin{equation}\label{eq:chern-character}
 \operatorname{ch}(\nabla)=\operatorname{tr}e^B
 =r+\sum_{k\geq1}\frac1{k!}\operatorname{tr}B^k.
\end{equation}
按底空间维数截断后仅有限项有贡献，故无分析收敛问题。对交换偶形式系数，Jacobi公式给
\[
 \frac{\dd}{\dd t}\log\det(I+tB)
 =\operatorname{tr}\bigl(B(I+tB)^{-1}\bigr).
\]
将$(I+tB)^{-1}$按形式幂级数展开，积分且用$t=0$常数项为零，得$\log\det(I+tB)=\operatorname{tr}\log(I+tB)$。比较系数得到
\begin{align*}
 c_1&=\operatorname{tr}B,&
 c_2&=\tfrac12\bigl((\operatorname{tr}B)^2-\operatorname{tr}B^2\bigr),\\
 \operatorname{ch}_1&=c_1,&
 \operatorname{ch}_2&=\tfrac12(c_1^2-2c_2),\\
 \operatorname{ch}_3&=\tfrac16(c_1^3-3c_1c_2+3c_3).
\end{align*}
这些首先是选定同一联络的形式恒等式，随后也是类恒等式。不能由$\operatorname{ch}$常带分母就否定整Chern类的归一化。

对直和联络，迹分块相加；对张量积联络，曲率为$F_E\otimes I+I\otimes F_F$。两项在偶形式代数中交换，指数可分解，且张量积矩阵的迹等于迹的乘积。因此
\[
 \operatorname{ch}(E\oplus F)=\operatorname{ch}(E)+\operatorname{ch}(F),\qquad
 \operatorname{ch}(E\otimes F)=\operatorname{ch}(E)\operatorname{ch}(F).
\]
这个证明不依赖尚未建立的拓扑分裂原理。对线丛$L$，$\operatorname{ch}(L)=e^{c_1(L)}$且$c_1(L\otimes K)=c_1(L)+c_1(K)$。

\subsection{实类与整类的界线}
第16--17章已经证明在$S^2$上圆群曲率周期的整数性，重言线丛满足$\int c_1(L)=-1$。但对任意秩、任意底空间，一般整类比较仍需拓扑构造及de Rham比较定理：存在拓扑整类$c_k^{\mathbb Z}(E)$，其到实上同调并与de Rham上同调识别后的像等于上述$c_k(E)$。\textbf{此一般比较在本书尚未证明。}后面的Pontryagin、Euler形式也采用相同层次的边界。即使承认比较，映到实系数后挠元消失；实类一般不能唯一还原整类。

\subsection{Pontryagin形式与稳定性}
实秩$n$向量丛$V$选内积及度量相容联络，曲率矩阵记$\Omega$，在正交标架中$\Omega^{\mathsf T}=-\Omega$。复化丛$V_{\mathbb C}$的曲率仍为同一矩阵。定义
\begin{equation}\label{eq:pontryagin}
 p_k(V)=(-1)^kc_{2k}(V_{\mathbb C}),\qquad
 p(\nabla;t)=\det\left(I+t\frac{\Omega}{2\pi}\right)
 =1+\sum_{k\geq1}t^{2k}p_k(\nabla).
\end{equation}
后式无奇数幂，因为转置给$\det(I+t\Omega)=\det(I-t\Omega)$；与Chern形式比较$i^{2k}=(-1)^k$即得前式。$p_k$为实$4k$次闭形式，其类不依赖联络、内积或定向，并满足Whitney乘法及加平凡实丛后的稳定性。特别地
\[
 p_1(\nabla)=-\frac1{8\pi^2}\operatorname{tr}(\Omega\wedge\Omega).
\]
若$E$为复丛，则$(E_{\mathbb R})_{\mathbb C}\simeq E\oplus\bar E$：把实复化中的乘$i$分成两特征丛，分别与$E$和$\bar E$识别即可。因此Whitney及共轭号给
\begin{equation}\label{eq:p1-underlying-complex}
 p_1(E_{\mathbb R})=c_1(E)^2-2c_2(E).
\end{equation}
这里实类的比较由丛同构及本章公式直接证明，不需整系数分裂原理。

\subsection{Pfaffian与定向Euler形式}
对实斜对称$2r$阶矩阵$X$定义
\[
 \operatorname{Pf}(X)=\frac1{2^r r!}
 \sum_{\sigma\in S_{2r}}\operatorname{sgn}(\sigma)
 X_{\sigma(1)\sigma(2)}\cdots X_{\sigma(2r-1)\sigma(2r)}.
\]
辅助外代数中$\Xi=\sum_{i,j}X_{ij}z_i\wedge z_j/2$满足$\Xi^r/r!=\operatorname{Pf}(X)z_1\wedge\cdots\wedge z_{2r}$。改变基并比较顶外幂系数，得到
\[
 \operatorname{Pf}(Q^{\mathsf T}XQ)=\det Q\,\operatorname{Pf}(X).
\]
所以它是$\mathrm{SO}(2r)$不变多项式，反向正交变换则改变符号。在实斜型的正交分块标准形上，Pfaffian是每个二阶块右上角元素的乘积，行列式是此乘积的平方；由换基律得到$\operatorname{Pf}(X)^2=\det X$。作为多项式恒等式，它也适用于偶形式系数。

\Needspace{18\baselineskip}
\begin{theorem}\label{thm:euler-form}
定向实秩$2r$丛$V$的Euler形式为
\begin{equation}\label{eq:euler-form}
 e(\nabla)=\frac{\operatorname{Pf}(\Omega)}{(2\pi)^r}.
\end{equation}
它闭，其类与度量相容联络及内积无关，反转丛定向使其变号。对定向偶秩丛有$e(V\oplus W)=e(V)e(W)$，且$p_r(V)=e(V)^2$。若$V$有无零光滑截面，则$e(V)=0$。
\end{theorem}
\begin{proof}
固定内积后应用$\mathrm{SO}(2r)$主丛的Chern--Weil构造及定理\ref{thm:cw-transgression}，得到闭性和联络无关性。反向标架的行列式为$-1$，由Pfaffian换基律给定向号。直和取按先$V$后$W$排列的定向标架，曲率分块；辅助外代数顶次项相乘给Euler乘法。Pfaffian平方等于行列式，再对照\eqref{eq:pontryagin}的最高项，得$p_r=e^2$。

若更换内积为$g_1(v,w)=g_0(Bv,w)$，$B$为$g_0$正定自伴场。正平方根光滑：在正根$C$处，映射$C\mapsto C^2$的导数为$T\mapsto CT+TC$，在对角基中各系数乘正数$\lambda_i+\lambda_j$，可逆，故逆函数定理给光滑依赖。于是$B^{-1/2}$为保向丛等距同构，将两个正交标架主丛识别；把第二个联络移回第一个主丛后才作仿射连接，得到同一Euler类。不能直接插值不同内积的相容联络，并默认为中间曲率仍在同一个$\mathfrak{so}$中。

最后把无零截面归一化，取分裂$V=\mathbb R\oplus V'$及使该截面平行的度量联络。其曲率首行首列为零，Pfaffian各项均含零因子，故此代表为零；联络无关性即给结论。
\end{proof}
复秩$r$丛按$(u_1,iu_1,\ldots,u_r,iu_r)$取底层实丛的复定向，有
\begin{equation}\label{eq:euler-top-chern}
 e(E_{\mathbb R})=c_r(E).
\end{equation}
证明可在多项式层完成：对$X\in\mathfrak u(r)$用酉对角化写成$i\operatorname{diag}(\lambda_1,\ldots,\lambda_r)$，其实表示的二阶块为$\left(\begin{smallmatrix}0&-\lambda_j\\\lambda_j&0\end{smallmatrix}\right)$。Pfaffian与$\det(iX)$均为$(-1)^r\prod_j\lambda_j$。酉换基的实表示保向，两边均不变，故对所有$X$相同；极化后代入曲率即得公式。这里仅用有限维Hermite谱定理，没有假设丛全局分裂为线丛。

\subsection{球面归一化及不消失向量场障碍}
第14章的二维Euler形式就是$K\mathrm{vol}/(2\pi)$。一般半径$R$的$S^{2r}$由第13章Gauss方程给常曲率$R^{-2}$，因而$\Omega_{ij}=R^{-2}\theta^i\wedge\theta^j$。Pfaffian的每个排列项都给同向体积形式，故
\[
 e(TS^{2r},\nabla)=
 \frac{(2r)!}{(2\pi)^r2^r r!\,R^{2r}}\mathrm{vol}.
\]
球面面积为$2\pi^{r+1/2}R^{2r}/\Gamma(r+1/2)$，可由Euclid高斯积分的径向形式求得；半整数Gamma递推给
\[
 \int_{S^{2r}}e=2.
\]
这是一组具体积分，不是一般Chern--Gauss--Bonnet的证明。结合无零截面必使Euler类为零，便证明偶维球面不存在处处非零切向量场。另一方面，$(X,a)\mapsto X+ap$给$TS^n\oplus\mathbb R\simeq S^n\times\mathbb R^{n+1}$，所以稳定性使所有正次Pontryagin类为零。Pontryagin类全零与Euler障碍可同时出现。


\Needspace{16\baselineskip}
\section{联络无关性与传递形式}
\subsection{两联络之间的全局相对形式}
设同一主丛上两联络为$\omega_0,\omega_1$。差$a=\omega_1-\omega_0$为水平$\operatorname{Ad}$等变一形式，因而是底空间的伴随丛值一形式。令$\omega_t=\omega_0+ta$，$0\leq t\leq1$，则
\begin{equation}\label{eq:curvature-interpolation}
 \Omega_t=\Omega_0+tD_0a+\tfrac{t^2}{2}[a,a],
 \qquad \dot\Omega_t=D_ta.
\end{equation}
这是在总空间展开$\dd\omega_t+[\omega_t,\omega_t]/2$得到的等式；矩阵情形最后一项为$t^2a\wedge a$。与$\omega_t$不同，$a$竖直收缩为零，因此$P(a,\Omega_t,\ldots)$也下降到底空间。

\Needspace{20\baselineskip}
\begin{theorem}[相对传递公式]\label{thm:cw-transgression}
对齐次$k\geq1$的不变$p$，定义全局$(2k-1)$形式
\begin{equation}\label{eq:cw-transgression}
 T_p(\omega_0,\omega_1)
 =k\int_0^1P(a,\Omega_t,\ldots,\Omega_t)\,\dd t.
\end{equation}
积分中的总空间形式均以下沉后的形式理解。则
\begin{equation}\label{eq:cw-difference}
 \operatorname{cw}_p(\omega_1)-\operatorname{cw}_p(\omega_0)
 =\dd T_p(\omega_0,\omega_1).
\end{equation}
因此Chern--Weil类与联络无关，且对丛同构及拉回自然。
\end{theorem}
\begin{proof}
对曲率的$k$个偶次输入求导，它们可无符号交换，故
\[
 \frac{\partial}{\partial t}P(\Omega_t,\ldots,\Omega_t)
 =kP(D_ta,\Omega_t,\ldots,\Omega_t).
\]
命题\ref{prop:invariant-derivative}用在首个一次输入$a$及其余二次输入上，再用$D_t\Omega_t=0$，得到右端等于$k\dd P(a,\Omega_t,\ldots,\Omega_t)$。系数对$t$光滑，在紧区间积分可逐坐标与外微分交换；积分即给所需差值。全部形式已下降，故这是$M$上的全局恰当性，不是仅在某张平凡化卡中的计算。自然性先在形式层成立，传递形式也与拉回交换，所以类的自然性不依赖选定联络。
\end{proof}
在$M\times[0,1]$上取无$\dd t$分量的联络$\omega_t$，其曲率为$\widetilde\Omega=\Omega_t+\dd t\wedge a$。代入$P$得到
\[
 P(\widetilde\Omega,\ldots,\widetilde\Omega)
 =P(\Omega_t,\ldots,\Omega_t)
 +k\dd t\wedge P(a,\Omega_t,\ldots,\Omega_t),
\]
右边省略的曲率输入补足到总共$k$个。对这个闭形式在参数方向积分，也得到差值公式；这解释了为何次数下降1。它不是把两底流形之间任意选一条路径来作平行移动。

\begin{figure}[htbp]
\centering\DGDiagram{connection-transgression}

\caption{联络仿射空间中的线段及参数积分。$T_p$为两联络的相对量，外微分等于两个特征形式之差；单个曲率代表随联络变化，而类保持不变。}
\label{fig:connection-transgression}
\end{figure}

\subsection{换连接路径与模恰当的可加性}
下文$P(a,\Omega^{k-1})$表示把$\Omega$重复填入$k-1$个位置，不是对$P$的值乘方。一般光滑联络路径也可用$k\int P(\dot\omega_t,\Omega_t^{k-1})\dd t$定义$T$。不同路径给出的$T$不必完全相等，但在固定端点时相差恰当形式。为看清这一点，取固定端点的两参数联络$\omega_{s,t}$，记$a=\partial_t\omega$、$b=\partial_s\omega$。直接微分并用Bianchi及$\partial_sa=\partial_tb$，得到
\begin{align*}
 &\partial_s\bigl(kP(a,\Omega^{k-1})\bigr)
 -\partial_t\bigl(kP(b,\Omega^{k-1})\bigr)\\
 &\hspace{1em}=-k(k-1)\dd P(a,b,\Omega^{k-2}).
\end{align*}
具体地，未消掉的项为
\[
 k(k-1)\bigl(P(a,D b,\ldots)-P(b,D a,\ldots)\bigr).
\]
另一方面，
\[
 \dd P(a,b,\ldots)=P(Da,b,\ldots)-P(a,Db,\ldots).
\]
$Da$与$b$交换无符号，故得到负号。固定端点使$t=0,1$处$b=0$，积分便得
\[
 T_{s=1}-T_{s=0}
 =-k(k-1)\dd\int_0^1\!\int_0^1P(a,b,\Omega^{k-2})\,\dd t\,\dd s.
\]
$k=1$时右端取零，直接积分$P(\dot\omega)$亦知结果只依赖端点。仿射空间中任意两条固定端点路径可线性同伦，因此结论适用。拼接路径的积分相加，遂有
\[
 T_p(\omega_0,\omega_2)\equiv
 T_p(\omega_0,\omega_1)+T_p(\omega_1,\omega_2)
 \pmod{\dd\Omega^{2k-2}(M)}.
\]
分段光滑拼接可逐段积分，或先作固定端点的平滑参数重标；没有将此模恰当关系写成未经选择的逐形式恒等式。

\subsection{局部Chern--Simons三形式}
在一个固定平凡化中，取$p(X)=\operatorname{tr}X^2$，以零势为起点、$A$为终点。$A_t=tA$的曲率为$t\dd A+t^2A\wedge A$，故相对传递式化为
\begin{equation}\label{eq:cs-three}
 \operatorname{CS}_3(A)=\operatorname{tr}
 \left(A\wedge\dd A+\tfrac23A\wedge A\wedge A\right),
 \qquad \dd\operatorname{CS}_3(A)=\operatorname{tr}(F\wedge F).
\end{equation}
也可直接核对：$\dd\operatorname{tr}(A\dd A)=\operatorname{tr}(\dd A\dd A)$，带次数循环律给$\dd\operatorname{tr}A^3=3\operatorname{tr}(\dd A A^2)$；另一方面$\operatorname{tr}A^4=0$，展开$\operatorname{tr}(\dd A+A^2)^2$正好得到同样结果。若用于$p_1$，还须整体乘$-1/(8\pi^2)$，不能省略归一化。

此处$A$是\emph{一张平凡化中的势}，零势也依赖该平凡化；所以$\operatorname{CS}_3(A)$一般不是非平凡丛上的全局基本形式。对比之下，式\eqref{eq:cw-transgression}固定了两个\emph{全局联络}，其差$a$为伴随丛值形式，因此给出全局相对形式。两者的全局性来源不同。

\subsection{纯规范势及二级信息的范围}
取$A=g^{-1}\dd g$，Maurer--Cartan式给$\dd A=-A^2$，所以$F=0$，但
\[
 \operatorname{CS}_3(g^{-1}\dd g)=-\tfrac13\operatorname{tr}(g^{-1}\dd g)^3
\]
不必为零。以$G=\mathrm{SU}(2)$为例，取$T_j=-i\sigma_j/2$，其中
\[
 \sigma_1=\begin{pmatrix}0&1\\1&0\end{pmatrix},\quad
 \sigma_2=\begin{pmatrix}0&-i\\i&0\end{pmatrix},\quad
 \sigma_3=\begin{pmatrix}1&0\\0&-1\end{pmatrix}.
\]
矩阵乘法给$[T_i,T_j]=\epsilon_{ijk}T_k$、$\operatorname{tr}(T_iT_jT_k)=-\epsilon_{ijk}/4$。写左Maurer--Cartan形式为$A=\sum_jT_j\theta^j$，逐项反对称求和给
\[
 \operatorname{tr}A^3=-\tfrac32\theta^1\wedge\theta^2\wedge\theta^3,
 \qquad \operatorname{CS}_3(A)=\tfrac12\theta^1\wedge\theta^2\wedge\theta^3\ne0.
\]
这已说明纯规范势下局部三形式可改变，虽然曲率为零。本章没有证明大规范变换使某个归一化积分只差整数，也没有构造微分上同调、平坦微分特征或规范不变的$\mathbb R/\mathbb Z$值泛函；不能仅由这个三形式便宣称所有二级理论已建立。

平坦联络使所有正次Chern--Weil类消失，因为直接取$F=0$作代表。但反向一般不成立，类消失也不等于丛平凡：第16章M\"obius线丛的$\mathrm O(1)$联络曲率恒零，丛仍非平凡，离散闭路移动为$-1$。这种结构的李代数为零，正次实不变多项式根本无法检出它。


\Needspace{16\baselineskip}
\section{示性数与怪球专题}
\subsection{闭流形上的示性数}
设$M$紧致、定向、无边界，$q$为若干Chern、Pontryagin或Euler形式的乘积，且总形式次数等于$\dim M$。定义相应示性数为$\int_Mq$。改变任一联络代表时，其差是恰当形式；与其余闭形式相乘仍恰当，Stokes使积分不变。拉回自然性及积分换元说明，保向丛同构覆盖的保向微分同胚保持此数，反转底定向则把一个\emph{固定}顶次形式积分变号。

此定义先得到实数。一般“整Chern/Pontryagin/Euler数为整数”的结论还需要前述整类比较及整基本类配对。我们将直接算出下面例子的整数值，不以少数例子代替一般整性定理。

\subsection{射影空间的积分与Euler短正合列}
在$\mathbb{CP}^m$上记$L$为重言线丛，$x=[\omega_{\mathrm{FS}}/(2\pi)]$。全纯标架$e_0=(1,z)$的内积系数为$1+|z|^2$，第17章Chern公式给$F_L=i\omega_{\mathrm{FS}}$，故$c_1(L)=-x$、$c_1(L^*)=x$。

我们先直接计算$\int_{\mathbb{CP}^m}x^m=1$。删去超平面后是$\mathbb C^m$，所删集合在有限坐标卡中实余维2，故对顶维体积积分为零。由第17章$\det G=(1+|z|^2)^{-(m+1)}$及体积公式，得到
\[
 \int x^m=\frac{m!}{\pi^m}
 \int_{\mathbb R^{2m}}\frac{\dd V}{(1+|z|^2)^{m+1}}.
\]
高斯积分给单位$S^{2m-1}$面积$2\pi^m/(m-1)!$；径向积分用$u=r^2$、$v=u/(1+u)$依次换元，得到
\[
 \int_0^\infty\frac{r^{2m-1}}{(1+r^2)^{m+1}}\dd r
 =\tfrac12\int_0^1v^{m-1}\dd v=\frac1{2m}.
\]
相乘即得1。此证明既确定尺度，也不预用$\mathbb{CP}^m$的整上同调环。

在直线$\ell\subset\mathbb C^{m+1}$处，图坐标把切向量识别为$\operatorname{Hom}_{\mathbb C}(\ell,\mathbb C^{m+1}/\ell)$：一条直线曲线的图函数一阶导数是一个从$\ell$到商空间的线性映射；换补空间只改变提升、不改变商。因此得到光滑且全纯相容的识别$T^{1,0}\mathbb{CP}^m\simeq\operatorname{Hom}(L,Q)$，$Q=\mathbb C^{m+1}/L$。对$0\to L\to\mathbb C^{m+1}\to Q\to0$逐纤维施加$\operatorname{Hom}(L,-)$，得
\begin{equation}\label{eq:projective-euler-sequence}
 0\longrightarrow\mathbb C\longrightarrow(L^*)^{\oplus(m+1)}
 \longrightarrow T^{1,0}\mathbb{CP}^m\longrightarrow0.
\end{equation}
第一项用线丛的恒等自同态识别为平凡线丛。第16章给出\emph{光滑}分裂，足以使用Whitney公式，无须声称存在全纯分裂。于是
\begin{equation}\label{eq:projective-chern}
 c(T^{1,0}\mathbb{CP}^m)=(1+x)^{m+1},\qquad
 c_k=\binom{m+1}{k}x^k\quad(0\leq k\leq m).
\end{equation}
在更高次数，底空间维数使相应de Rham类为零；没有把形式上存在的$x^{m+1}$当成一个额外非零类。

\subsection{四维计算与有边界的区别}
复定向$\mathbb{CP}^2$上，公式给
\[
 \int c_1^2=9,\qquad \int c_2=3,\qquad
 \int p_1=\int(c_1^2-2c_2)=3,\qquad \int e=3.
\]
$\mathbb{CP}^1\times\mathbb{CP}^1$上，令$x,y$为两因子归一化面积形式的类，$x^2=y^2=0$、$\int xy=1$由维数及乘积积分直接给出。切丛分成两因子，故$c_1=2x+2y$、$c_2=4xy$，从而$\int c_1^2=8$、$\int c_2=\int e=4$、$\int p_1=0$。Pontryagin数已足以区分这两个光滑流形，即使允许反向微分同胚，$\pm3$也不等于0。

若$M=\partial W$且给定丛及闭特征形式确实延伸到紧致定向$W$，外向首位边界定向及Stokes给$\int_Mq=\int_W\dd q=0$。对切丛Pontryagin数，可用$TW|_M\simeq TM\oplus\mathbb R$及稳定性，故它们是相应定向配边关系下的不变量。不能将这个论证照搬到切丛Euler数：$TW$的秩及其Euler形式次数不同，Euler类也不稳定；$S^2=\partial D^3$却有$\int_{S^2}e=2$，就是直接反例。

\begin{figure}[htbp]
\centering\DGDiagram{characteristic-limits}

\caption{前两例取复定向。具体Pontryagin数可以区分某些四维光滑流形，但球面型七维流形的正次实示性类消失仍不能给出完整光滑分类。各框表示检测结果，未宣称后者与标准球面微分同胚。}
\label{fig:characteristic-limits}
\end{figure}

\subsection{为什么七维球面需要别的信息}
称光滑闭流形为怪球（exotic sphere），是指它与标准球面同胚但不微分同胚。“同胚”和“微分同胚”不应只因两者都有球面外形而合并。Milnor于1956年构造了七维这样的例子；这是外部定理，本章只说明Chern--Weil初等检测在此处的局限，不声称完成怪球构造或分类。

若一个七维流形在光滑同伦意义下与$S^7$等价，第9章使其de Rham上同调仅在次数0、7非零。正次Pontryagin类的次数为$4,8,\ldots$，故全部为零；一般Chern--Weil偶次类也无法在次数7给出顶次示性数。对只给同胚的怪球使用同样的上同调结论时，还须采用本书未证明的拓扑/de Rham比较与相关光滑逼近，不把这一比较隐藏在“同胚不变”四字中。

这并不否定Pontryagin理论在Milnor证明中的作用。区别在于，他使用八维填充及相对配对，而不是仅积分七维流形自身的四形式。为准确说明这一点，给定紧致定向八维$W$，$\partial W=M^7$，并假设整上同调$H^3(M;\mathbb Z)=H^4(M;\mathbb Z)=0$。相对上同调长正合列使
\[
 j:H^4(W,M;\mathbb Z)\longrightarrow H^4(W;\mathbb Z)
\]
为同构。\textbf{以下属于外引拓扑工具：}取整Pontryagin类$p_1^{\mathbb Z}(TW)$，定义
\[
 q(W)=\left\langle\bigl(j^{-1}p_1^{\mathbb Z}(TW)\bigr)^2,[W,M]\right\rangle,
 \qquad\lambda(M)=2q(W)-\sigma(W)\pmod7,
\]
其中$\sigma(W)$是相对四维交叉配对的签名。Milnor原论文定理1断言此剩余类与填充的选择无关，并以它判别某些球面束构造所得流形的光滑结构。该论证使用整类、相对基本类、填充存在性和Hirzebruch签名定理；这些并没有由本章实形式的联络无关性自动推出。

因此不能用$H^4(M;\mathbb R)=0$断言$q(W)=0$，也不能把某张卡上的Chern--Simons三形式随意积分后叫作上述$\lambda$。本章不证明一般Chern--Gauss--Bonnet、Euler类与零点指标比较、整类比较、Hirzebruch签名公式、怪球分类或高维手术理论。第14章已经证明的曲面公式和本章算出的球面/射影积分，均不冒充这些整体定理的一般证明。


\Needspace{18\baselineskip}
\section{示性类习题及详细解答}
18.1--18.3分别以Freed《Geometry of Dirac Operators》约1987年初稿的题6.1.13、6.1.15、6.1.16为起点改述扩展，均位于印刷92/PDF94；本书改用$iF/(2\pi)$并纠正该页括号中超交换子号及幂次的排印问题。18.4--18.8为原创补充。一般定理已在正文证明，这些习题不承担缺失的主线证明。

\Needspace{12\baselineskip}
\subsection*{习题18.1：有符号循环律与反例}
改述Freed题6.1.13并扩展。证明矩阵值齐次形式的带次数循环律及$\dd\operatorname{tr}\alpha=\operatorname{tr}D_A\alpha$。取$\alpha=E_{12}\dd x$、$\beta=E_{21}\dd y$，比较普通交换子与分次交换子的迹。
\begin{proof}[解答]
按指标展开$\operatorname{tr}(\alpha\beta)=\sum_{i,j}\alpha_{ij}\wedge\beta_{ji}$，移动标量形式并互换$i,j$，得$(-1)^{rs}\operatorname{tr}(\beta\alpha)$。所以分次交换子必须取$[\alpha,\beta]_{\mathrm{gr}}=\alpha\beta-(-1)^{rs}\beta\alpha$，其迹为零。$D_A\alpha$与$\dd\alpha$之差是同样的分次交换子，得到求导公式。

本例$\operatorname{tr}(\alpha\beta)=\dd x\wedge\dd y$、$\operatorname{tr}(\beta\alpha)=-\dd x\wedge\dd y$，故普通交换子的迹为$2\dd x\wedge\dd y$，分次交换子此时为反交换子，迹为零。源页括号中若取印刷的加号约定，在零次矩阵情形便已经与“迹为零”冲突；本书按上述直接计算固定减号。
\end{proof}

\Needspace{12\baselineskip}
\subsection*{习题18.2：闭迹幂及Chern特征的展开}
改述Freed题6.1.15并扩展。证明$\operatorname{tr}F^k$闭，写出$\operatorname{ch}$至六次的展开，并求直和与张量积公式。说明指数为何可视为有限表达式。
\begin{proof}[解答]
协变外微分按分次乘法式作用，$F$为偶次，故$D_AF^k=\sum_jF^j(D_AF)F^{k-1-j}=0$。再用上一题的迹求导公式得到闭性。底空间维数有限，高于维数的形式均为零，所以指数在形式意义下截断，不需估计矩阵指数的分析收敛。

记$B=iF/(2\pi)$，逐阶展开$\operatorname{tr}e^B$，用正文的行列式恒等式消去迹幂，得
\[
 \operatorname{ch}=r+c_1+\tfrac12(c_1^2-2c_2)
 +\tfrac16(c_1^3-3c_1c_2+3c_3)+\text{高于六次的项}.
\]
直和的指数按块分解，迹相加。张量积曲率两项$B_E\otimes I$和$I\otimes B_F$交换，指数积的迹为两迹之积，得到$\operatorname{ch}(E\otimes F)=\operatorname{ch}(E)\operatorname{ch}(F)$。这是对源题闭性主题的扩展，不把新增的张量积结论冒充原题原句。
\end{proof}

\Needspace{12\baselineskip}
\subsection*{习题18.3：指数传递及三次系数}
改述Freed题6.1.16并按本书归一化扩展。若$\dot A_t=a_t$，证明
\[
 \frac{\partial}{\partial t}\operatorname{tr}e^{\lambda F_t}
 =\lambda\dd\operatorname{tr}(a_t e^{\lambda F_t}),
\]
其中$\lambda$为常数；再由线段$A_t=tA$求$\operatorname{tr}F^2$的三次传递形式。它是否自动在非平凡丛上全局定义？
\begin{proof}[解答]
不交换矩阵的指数不能直接按标量链式法则求导，但对每一迹幂有
\[
 \partial_t\operatorname{tr}F_t^k
 =k\operatorname{tr}(\dot F_tF_t^{k-1})
 =k\dd\operatorname{tr}(a_tF_t^{k-1}),
\]
第一步的$k$项通过偶次循环律变为相同，第二步用$\dot F_t=D_ta_t$和Bianchi。将有限截断指数逐阶求和便得公式。取$\lambda=i/(2\pi)$给本书Chern特征，源书则取$\lambda=-1$。

线段曲率为$t\dd A+t^2A^2$，所以
\[
 2\int_0^1\operatorname{tr}(A F_t)\dd t
 =\operatorname{tr}(A\dd A+\tfrac23A^3).
\]
这是固定平凡化中的$\operatorname{CS}_3(A)$。非平凡丛上“零势”未必拼成一个全局联络，因此此式不自动全局。若确有两个全局联络，改用它们的差及相对传递才有全局性。
\end{proof}

\Needspace{12\baselineskip}
\subsection*{习题18.4：射影平面上线丛直和}
原创补充。令$L\to\mathbb{CP}^2$为重言线丛，$E=L^a\oplus L^b$，$a,b\in\mathbb Z$，负幂表示对偶幂。以$x=[\omega_{\mathrm{FS}}/(2\pi)]$写$c_1,c_2,\operatorname{ch}_2,p_1(E_{\mathbb R})$，并求对应四次积分。
\begin{proof}[解答]
线丛张量及对偶公式给$c_1(L^a)=-ax$，所以
\begin{align*}
 c_1(E)&=-(a+b)x,& c_2(E)&=abx^2,\\
 \operatorname{ch}_2(E)&=\tfrac12(a^2+b^2)x^2,&
 p_1(E_{\mathbb R})&=(a^2+b^2)x^2.
\end{align*}
由于$\int x^2=1$，$c_1^2,c_2,\operatorname{ch}_2,p_1$的积分依次为$(a+b)^2,ab,(a^2+b^2)/2,a^2+b^2$。$\operatorname{ch}_2$可为半整数，这不违反Chern数的整性归一化。若取$E=L$本身，则底层实二平面丛有$e=-x$、$p_1=x^2=e^2$。
\end{proof}

\Needspace{12\baselineskip}
\subsection*{习题18.5：Pfaffian的定向与复表示号}
原创补充。在有序实基$(u,iu,v,iv)$中取$X=\operatorname{diag}(ia,ib)\in\mathfrak u(2)$的实表示，计算Pfaffian及行列式；交换前两个基向量后会怎样？用它核对$e(E_{\mathbb R})=c_2(E)$与$p_2=e^2$的号。
\begin{proof}[解答]
实表示为两个块$\left(\begin{smallmatrix}0&-a\\a&0\end{smallmatrix}\right)$、$\left(\begin{smallmatrix}0&-b\\b&0\end{smallmatrix}\right)$，所以$\operatorname{Pf}=ab$，行列式$a^2b^2$。同时$\det(iX)=(-a)(-b)=ab$，与复定向的顶Chern多项式一致。交换前两个基向量使换基行列式为$-1$，Pfaffian变成$-ab$，行列式仍为$a^2b^2$。

因此Euler形式需要丛定向，Pontryagin形式不需要；顶Pontryagin只看到Euler的平方。底流形定向则另控制积分号。对切丛同时反转丛定向和底定向时，Euler形式及积分方向各变一次号，Euler数保持不变，不可把这两种定向变化合并成一次。
\end{proof}

\Needspace{12\baselineskip}
\subsection*{习题18.6：相对一阶传递与曲率零}
原创补充。同一Hermite线丛上两个酉联络相差$i\alpha$，$\alpha$为全局实一形式。计算$c_1$形式之差及其传递形式。在平凡线丛上取$A_0=0,A_1=ix\dd y$；再以M\"obius线丛说明零曲率为什么不等于平凡。
\begin{proof}[解答]
圆群交换，故$F_1-F_0=i\dd\alpha$，乘$i/(2\pi)$得到
\[
 c_1(\nabla_1)-c_1(\nabla_0)
 =-\frac{\dd\alpha}{2\pi},\qquad T_{c_1}=-\frac\alpha{2\pi}.
\]
给定例子$F_1=i\dd x\wedge\dd y$，$c_1(\nabla_1)=-\dd x\wedge\dd y/(2\pi)$非零，但它为全局恰当形式，类仍为零。闭曲面上的积分不会因此改变。

M\"obius实线丛的单位标架丛结构群为$\mathrm O(1)$，李代数为零，唯一联络曲率也为零，所有正次Pontryagin类消失。非平凡性却由绕圈反号或无零截面障碍给出。本例涉及实线丛，并非声称已构造某个任意复秩的全部Chern类零但非平凡的分类反例。
\end{proof}

\Needspace{12\baselineskip}
\subsection*{习题18.7：乘积曲面的两组示性数}
原创补充。对$M=\mathbb{CP}^1\times\mathbb{CP}^1$取乘积复结构，计算$c_1^2,c_2,p_1,e$的积分，并与$\mathbb{CP}^2$比较。哪一部分推理依赖一般Chern--Gauss--Bonnet？
\begin{proof}[解答]
写$x,y$为两因子归一化面积类，$x^2=y^2=0$、$\int xy=1$。每个射影直线切丛有$c_1=2x$或$2y$，切丛直和使
\begin{gather*}
 c(TM)=(1+2x)(1+2y),\\
 c_1^2=8xy,\quad c_2=4xy,\quad p_1=0,\quad e=c_2.
\end{gather*}
积分依次为$8,4,0,4$。射影平面的对应结果为$9,3,3,3$，所以Pontryagin数已排除两者微分同胚，允许反向也不改变“非零与零”的区别。这里Euler等于顶Chern由酉表示的多项式恒等式证明，所有积分由显式FS形式或乘积积分得到；并未用一般Euler数等于拓扑Euler示性数。若将这些数统一解释为任意闭流形的$\chi$，那一步才需外部比较或Chern--Gauss--Bonnet。
\end{proof}

\Needspace{12\baselineskip}
\subsection*{习题18.8：稳定平凡性与七维检测的局限}
原创补充。显式证明$TS^n\oplus\mathbb R$平凡。为什么这会使正次Pontryagin类全零，却不使偶维球面Euler类为零？对一个光滑同伦七维球面，Chern--Weil正次偶类能否区分其与标准球面的光滑结构？
\begin{proof}[解答]
在单位球面点$p$处，把$(X,a)$送为$X+ap\in\mathbb R^{n+1}$。逆为$v\mapsto(v-\langle v,p\rangle p,\langle v,p\rangle)$，两式光滑且逐纤维线性互逆。Pontryagin类稳定且平凡丛的正次类为零，故$p(TS^n)=1$。Euler类却不稳定；本章圆球直接积分给$\int_{S^{2r}}e=2$，所以其类非零。不能通过加一条平凡线再把奇秩曲率的行列式公式当作原Euler类。

光滑同伦七维球面的正偶次数de Rham群均为零，且奇数顶维7无法由偶次数相加得到，所以这些类及其普通示性数不能完成区分。这不是微分同胚结论。Milnor的外部论证使用八维填充、整相对配对及签名；本题没有计算这些二级数据，亦未证明任何怪球分类总数。
\end{proof}

